   
\section{Mobility patterns and user behavior}
% \label{sec:mobility_patterns_and_user_behaviour}

% Literally copied from Zhang2018 pg 220
% After retrieving the chronologically ordered trips, we calculated the distances of all trips. We deleted those with a distance longer than 50 km or a time duration longer than 10 h. After this process of data cleaning, we obtained 1,023,529 trips. The estimated average trip distance was 2.4 km, with a standard deviation of 2.2 km. The total distance of all bike trips was 2.4 million km. The average travel time was 16.8 min, with a standard deviation of 18.5 min. The travel time for all bike trips was equivalent to 32.7 years. The spatial distribution of all bike trips (shown in Fig. 2) was created using the line density method in

% Copiat literal de Bonilla-Alicea2019 pg 139
% Evidence indicates, however, that smart bike BSS users may have less sustainable usage patterns. In the United States, smart bike systems averaged 0.3 rides per bike per day, while smart dock systems averaged 1.7 rides per bike per day in 2017 (National Association of City Transportation Officials, 2017). Other smart bike ridership estimates range from 0.5 to 1.6 rides per bike per day (Mynorthwest.com, 2018; Relay Bike Share, 2017; Runyan, 2017). In Seattle, smart dock systems demonstrate usage patterns that correlate with commuting rush hours, while smart bike systems exhibit trip patterns suggestive of recreational use (National Association of City Transportation Officials, 2017). InWashington, DC, smart bikes have more geographically diverse usage, most likely because trips are not constrained to docking stations (Virginia Tech, 2018). Despite fulfilling different functional needs, both smart dock and smart bike BSS trips inWashington, DC are generally less than 3 miles.

% Julio2022 Long term assessment of a successful e-bike-sharing system. Key drivers and impact on travel behaviour

%Ricci2015

% Fisherman
% Trip durations, purposes, distances, demographics, etc. .

% Temporal patterns (rush hour vs leisure).

% One-way vs round-trip.

% User typologies (subscribers vs casual users).

% First–last mile use cases and multimodal integration.

% % ###################################################
% % https://www.researchgate.net/publication/264356705_Understanding_the_Diffusion_of_Public_Bike-sharing_Systems_Evidence_from_Europe_and_North_America
% % The principle of bike-sharing systems is simple: bike-sharing users access bicycles on an as-needed basis. Public bike-sharing stations are typically unattended and concentrated in urban settings. They provide a variety of pickup and drop-off locations, enabling an on-demand, very low emission form of mobility. The majority of bike-sharing programs cover the costs of bicycle maintenance, storage, and parking (similar to carsharing or short-term auto access). Trips can be  point-to-point, round-trip, or both, allowing the bikes to be used  for one-way transport and for multimodal connectivity (first-and-last mile trips, many-mile trips, or both) (Shaheen et al.,  2013; Shaheen et al., 2012a). The last mile refers to  the distance between workplaces or homes and the public transport stops where users have disembarked (Shaheen et al., 2010). If these distances are too great to walk in a reasonable time, bike-sharing offers users an option to help them complete their journey.

% % Generally, trips of less than 30 minutes are covered through a daily, monthly, and annual pass at no extra charge. They can pick up a bike at any dock by using their credit or debit card, membership card, or key, and/or a mobile phone. When they finish using the bike, they can return it to any dock (or the same dock in a round-trip service) where there is a spot and end their  session.  By  addressing  the  storage,  maintenance,  and  parking  aspects  of  bicycle ownership, public bike-sharing encourages cycling among users who may not otherwise ride bikes. Additionally, the availability of a large number of bicycles in multiple dense, nearby locations  frequently  creates  a  “network-effect,”  further  encouraging  cycling  and,  more specifically,  the  use  of  public  bike-sharing  for  regular  trips  (e.g.,  commuting,  errands) (Shaheen et al., 2012a).